\documentclass[11pt]{article}
\usepackage{amsmath,amssymb}
\title{Local vector, gauge-Jacobian and Haar vertices for the first nonlinear coefficient}\author{}\date{}
\begin{document}\maketitle
\paragraph{Why the quadratic quotient is insufficient at the next order.}
The exact tree covariance can be written $T\Delta^{-1}T^*$, but replacing
the propagator alone inside a tree-coordinate quartic expression does
not produce a local covariant formula. The nonlinear gauge-coordinate
Jacobian has vertices. This note retains them and verifies the resulting
first coefficient on the original complete $R=2$ cube.

Fix $R$ and the admitted actual action minimum $U^*(\epsilon)$ on all
free links. Interior gauge parameters vanish on the complete boundary.
Use right logarithmic coordinates $U_e=U_e^*\exp(\alpha z_e)$,
$\alpha=B^{-1/2}$, with the $SU(2)$ convention $z=i{\bf z}\cdot\sigma$.
Define the two local endpoint maps
\[
 (D\xi)_e=\operatorname{Ad}_{U_e^{*-1}}\xi_{s(e)}-\xi_{t(e)},\qquad
 (A\xi)_e=\operatorname{Ad}_{U_e^{*-1}}\xi_{s(e)}+\xi_{t(e)}.
\]
The relative gauge matrix $\Delta_0=D^*D$ is positive, and the actual
critical action Hessian $H$ obeys $HD=0$. The full vector covariance is
$\Delta_1^{-1}$, $\Delta_1=H+DD^*$.

\paragraph{Literal local gauge-Jacobian expansion.}
Choose the local logarithmic gauge function $F(z)=D^*z$, with $D$
held at the current minimum. Under
$U_e\mapsto e^{\alpha t\xi_s}U_e e^{-\alpha t\xi_t}$, the BCH derivative is
\[
 \left.\partial_t z_e\right|_0
 = (D\xi)_e-\frac{\alpha}{2}[z_e,(A\xi)_e]
       +\frac{\alpha^2}{12}[z_e,[z_e,(D\xi)_e]]+O(\alpha^3).
\]
Consequently its actual Jacobian in the interior gauge parameter is
\[
 M(z)=\Delta_0+\alpha M_1(z)+\alpha^2 M_2(z)+\cdots,\quad
 M_1=-\tfrac12D^*\operatorname{ad}_z A,\quad
 M_2=\tfrac1{12}D^*\operatorname{ad}_z^2 D .
\]
Let $g=\Delta_0^{-1}$. Its logarithmic determinant partners are
\[
 \ell_1=\operatorname{Tr}(gM_1),\qquad
 \ell_2=\operatorname{Tr}(gM_2)-\tfrac12\operatorname{Tr}(gM_1gM_1).
                                                               \tag{1}
\]
These contain the original local gauge vertices and their propagator;
they are not a constant extra determinant. At the flat reference,
$\ell_1\equiv0$ by the color trace of $\operatorname{ad}_z$, but its
background derivatives must still be retained.

\paragraph{Precise domain of the gauge statement.}
This construction is LOCAL about the minimum orbit at each fixed $R$.
The inverse function theorem for the relative matrix $\Delta_0$ supplies
a unique gauge intersection for sufficiently small coordinates and gauge
noise. One can insert the local Jacobian and a normalized Gaussian
integral over the gauge function, with a smooth partition equal to one
near the minimum orbit. The resulting quadratic gauge penalty is
$|D^*z|^2/2$. Outside that neighborhood the original tree-fixed compact
integral has the positive fixed-$R$ action gap used in the first original nonlinear flux coefficient calculation.
Thus its Laplace coefficients agree with these local gauge-fixed
coefficients. The local construction is not an exact global
Faddeev--Popov formula: no assertion about distant gauge copies,
uniform patch size, or a growing-$R$ complementary integral is made.
Gauge orbit volume and all constant canonical factors are retained.
At quadratic order its normalization is
$\det\Delta_0/(\det\Delta_1)^{1/2}$, equal to the original tree factor.

\paragraph{The complete local first coefficient.}
Let $T_3,T_4$ be the ORIGINAL full-free-link action Taylor vertices in
these right coordinates, and
$J_2=-\sum_{\text{free }e}|{\bf z}_e|^2/3$ the full Haar polynomial.
Expansion of the actual action, Haar density and Jacobian gives the
same finite-block coefficient as the tree calculation:
\[
 \boxed{\quad
 \mathcal C_R(\epsilon)=
 E_{\Delta_1^{-1}}\left[
 T_4-J_2-\ell_2-\tfrac12(T_3-\ell_1)^2\right].\quad}              \tag{2}
\]
Indeed the order-$\alpha$ exponent is $-T_3+\ell_1$, and its square,
the order-$\alpha^2$ exponent $-T_4+\ell_2$, and Haar $J_2$ supply the
order-$B^{-1}$ coefficient of the integral. Taking minus its logarithm
gives (2); the odd first term has zero centered Gaussian expectation.
This equality concerns the complete combination, not a term-by-term
identification with tree vertices. At a varying background all
derivatives of $D,A,g,\Delta_1$ and the original action tensors are kept.
The normalized connected differentiation formula of the connected first nonlinear flux response calculation applies
with the complete bracket in (2).

\paragraph{Exact original $R=2$ validation.}
At flat $R=2$ there is one interior vertex, eight free links and
$\Delta_0=8I_3$. Write $d_e=\pm1$ for its signed incidence.
Each plaquette has at most two free links, so $T_3=0$.
The scalar vector Green matrix consists of four two-by-two blocks
with diagonal $7$ and off-diagonal $-1$ before inversion. Hence
$\operatorname{Tr}G=7/6$ and $d^TGd=1$.
Using $[i{\bf z}\cdot\sigma,i{\bf \xi}\cdot\sigma]
=-2i({\bf z}\times{\bf \xi})\cdot\sigma$ gives
\[
 M_1=\operatorname{cross}\left(\sum_e d_e{\bf z}_e\right),\qquad
 M_2=\tfrac13\sum_e\operatorname{cross}({\bf z}_e)^2.
\]
Original plaquette quartic products and exact Gaussian contraction yield
\[
 ET_4=-\frac{343}{192},\qquad -EJ_2=\frac76,\qquad
 E\ell_2=-\frac14\operatorname{Tr}G+\frac3{64}d^TGd
                  =-\frac{47}{192}.
\]
Thus (2) gives
\[
 \mathcal C_2(0)=-\frac{343}{192}+\frac{224}{192}
                         +\frac{47}{192}=-\frac38,
\]
exactly matching the independent tree calculation. The finite calculation verifies every term.
Omitting the nonlinear gauge Jacobian would fail this test.

\paragraph{Uniformity target, now with explicit partners.}
Formula (2) supplies local action and gauge vertices with full relative
vector and scalar Green operators, rather than long axial-tree rows.
A physical-energy estimate must bound its twice differentiated CONNECTED
combination, including the Haar term and both terms in (1).
The quadratic magnetic-loop bound does not estimate these interacting
multi-vertex contractions or the original differentiated remainder.
No bound uniform in $R$, random-background extension, or S94 closure is
claimed. The local gauge identity makes the missing estimate explicit;
it does not supply that estimate by itself.
\end{document}
